\documentclass[12pt]{article}
\usepackage[T2A]{fontenc}
\usepackage[utf8]{inputenc}
\usepackage[russian]{babel}
\usepackage{amsmath, amssymb}
\usepackage{geometry}
\usepackage{hyperref}
\usepackage{verbatim}
\geometry{margin=1.2cm}

\title{Решётка по проверочной матрице (hom6)}
\author{}
\date{}

\begin{document}
\maketitle

\section*{Постановка}
Порождающая матрица кода:
\[
G=
\begin{pmatrix}
1 & 0 & 1 & 1 & 0 & 1\\
1 & 0 & 1 & 0 & 1 & 0\\
1 & 1 & 0 & 1 & 0 & 0
\end{pmatrix}.
\]

Требуется:
1) построить решётку по проверочной (проверочной) матрице;
2) убедиться, что полученная решётка совпадает (с точностью до переупорядочивания узлов на ярусах) с синдромной решёткой этого же кода.

\section*{Идея построения}
Код задаётся линейным условием
\[
H x^T = 0 \quad (\text{над } \mathbb{F}_2),
\]
где $H$ --- матрица проверок.

Для решётки (синдромной/по $H$) на ярусе $t$ используется \emph{частичный синдром}
\[
s_t = H_{[:,0..t-1]} \, x_{0..t-1}^T,
\]
а переход при выборе следующего бита $b=x_t$ задаётся правилом
\[
s_{t+1} = s_t + b \cdot H_{[:,t]}.
\]
При этом узлы берутся только те, которые достижимы некоторыми кодовыми словами.

\section*{Что делает программа}
Файл \texttt{main.py}:
1) находит матрицу $H$ как базис правого нуля матрицы $G$ (то есть набор векторов $h$, для которых $G h^T=0$);
2) перечисляет все кодовые слова;
3) строит две решётки:
   \begin{itemize}
   \item \texttt{Т1: решетка по H} (строится из проверочных уравнений с использованием частичных синдромов и достижимости);
   \item \texttt{Т2: синдромная решетка} (строится напрямую: по частичным синдромам каждого кодового слова, затем рёбра определяются переходами между ярусами).
   \end{itemize}
4) проверяет совпадение решёток программно: узлы сравниваются по тому, \emph{какие кодовые слова} проходят через них (то есть совпадает структура решётки, а переупорядочивание узлов на уровнях допускается).

\section*{Вывод программы}
Итоговая проверка совпадения решёток (фрагмент из \texttt{out.txt}):

\begin{verbatim}
   G (k x n) =
[1 0 1 1 0 1]
[1 0 1 0 1 0]
[1 1 0 1 0 0]
rank(G) = 3, k=3, n=6

Проверочная матрица H (как базис nullspace(G)):
[1 1 1 0 0 0]
[1 0 0 1 1 0]
[0 1 0 1 0 1]
размер H: 3 x 6

Проверка G * H^T mod 2 (должно быть нулевой матрицей):
[0 0 0]
[0 0 0]
[0 0 0]

Всего различных кодовых слов: 8
c[0] = [0, 0, 0, 0, 0, 0]
c[1] = [1, 1, 0, 1, 0, 0]
c[2] = [1, 0, 1, 0, 1, 0]
c[3] = [0, 1, 1, 1, 1, 0]
c[4] = [1, 0, 1, 1, 0, 1]
c[5] = [0, 1, 1, 0, 0, 1]
c[6] = [0, 0, 0, 1, 1, 1]
c[7] = [1, 1, 0, 0, 1, 1]

===== Решетка: Т1: решетка по H =====
H (размер 3x6):
   1 1 1 0 0 0
   1 0 0 1 1 0
   0 1 0 1 0 1

Уровень t=0, число узлов: 1
  node 000: codewords=[0, 1, 2, 3, 4, 5, 6, 7]
  Рёбра:
    000 --(b=0)-> 000
    000 --(b=1)-> 110

Уровень t=1, число узлов: 2
  node 000: codewords=[0, 3, 5, 6]
  node 110: codewords=[1, 2, 4, 7]
  Рёбра:
    000 --(b=0)-> 000
    000 --(b=1)-> 101
    110 --(b=0)-> 110
    110 --(b=1)-> 011

Уровень t=2, число узлов: 4
  node 000: codewords=[0, 6]
  node 011: codewords=[1, 7]
  node 101: codewords=[3, 5]
  node 110: codewords=[2, 4]
  Рёбра:
    000 --(b=0)-> 000
    011 --(b=0)-> 011
    101 --(b=1)-> 001
    110 --(b=1)-> 010

Уровень t=3, число узлов: 4
  node 000: codewords=[0, 6]
  node 001: codewords=[3, 5]
  node 010: codewords=[2, 4]
  node 011: codewords=[1, 7]
  Рёбра:
    000 --(b=0)-> 000
    000 --(b=1)-> 011
    001 --(b=0)-> 001
    001 --(b=1)-> 010
    010 --(b=0)-> 010
    010 --(b=1)-> 001
    011 --(b=0)-> 011
    011 --(b=1)-> 000

Уровень t=4, число узлов: 4
  node 000: codewords=[0, 1]
  node 001: codewords=[4, 5]
  node 010: codewords=[2, 3]
  node 011: codewords=[6, 7]
  Рёбра:
    000 --(b=0)-> 000
    001 --(b=0)-> 001
    010 --(b=1)-> 000
    011 --(b=1)-> 001

Уровень t=5, число узлов: 2
  node 000: codewords=[0, 1, 2, 3]
  node 001: codewords=[4, 5, 6, 7]
  Рёбра:
    000 --(b=0)-> 000
    001 --(b=1)-> 000

Уровень t=6, число узлов: 1
  node 000: codewords=[0, 1, 2, 3, 4, 5, 6, 7]

===== Решетка: Т2: синдромная решетка (из перечисления кодовых слов) =====
H (размер 3x6):
   1 1 1 0 0 0
   1 0 0 1 1 0
   0 1 0 1 0 1

Уровень t=0, число узлов: 1
  node 000: codewords=[0, 1, 2, 3, 4, 5, 6, 7]
  Рёбра:
    000 --(b=0)-> 000
    000 --(b=1)-> 110

Уровень t=1, число узлов: 2
  node 000: codewords=[0, 3, 5, 6]
  node 110: codewords=[1, 2, 4, 7]
  Рёбра:
    000 --(b=0)-> 000
    000 --(b=1)-> 101
    110 --(b=0)-> 110
    110 --(b=1)-> 011

Уровень t=2, число узлов: 4
  node 000: codewords=[0, 6]
  node 011: codewords=[1, 7]
  node 101: codewords=[3, 5]
  node 110: codewords=[2, 4]
  Рёбра:
    000 --(b=0)-> 000
    011 --(b=0)-> 011
    101 --(b=1)-> 001
    110 --(b=1)-> 010

Уровень t=3, число узлов: 4
  node 000: codewords=[0, 6]
  node 001: codewords=[3, 5]
  node 010: codewords=[2, 4]
  node 011: codewords=[1, 7]
  Рёбра:
    000 --(b=0)-> 000
    000 --(b=1)-> 011
    001 --(b=0)-> 001
    001 --(b=1)-> 010
    010 --(b=0)-> 010
    010 --(b=1)-> 001
    011 --(b=0)-> 011
    011 --(b=1)-> 000

Уровень t=4, число узлов: 4
  node 000: codewords=[0, 1]
  node 001: codewords=[4, 5]
  node 010: codewords=[2, 3]
  node 011: codewords=[6, 7]
  Рёбра:
    000 --(b=0)-> 000
    001 --(b=0)-> 001
    010 --(b=1)-> 000
    011 --(b=1)-> 001

Уровень t=5, число узлов: 2
  node 000: codewords=[0, 1, 2, 3]
  node 001: codewords=[4, 5, 6, 7]
  Рёбра:
    000 --(b=0)-> 000
    001 --(b=1)-> 000

Уровень t=6, число узлов: 1
  node 000: codewords=[0, 1, 2, 3, 4, 5, 6, 7]

===== Сравнение решёток =====
Совпадение: True
Причина: Требуемое совпадение выполнено (с точностью до переупорядочивания узлов на уровнях).

\end{verbatim}

\end{document}

